\section{Fixed Methods}\label{sec:fixed}

\subsection{Weak and strong linearisation-dependent solutions}

For the weak and strong LD columns the published results are correct, and we only sketch the
arguments, pointing out where they had to be adjusted.

\paragraph{Weak solutions.}
For primitive problems, weak plan existence is \NP-complete in both the totally and the
partially ordered case \parencite[Thm.~4.1, Cor.~4.2]{chen:bercher:icaps-21}. Membership is
guess-and-check. Hardness comes from a reduction from SAT in which each variable is an action with
two outcomes, adding the clauses that the true resp.\ false literal satisfies, followed by a task
that requires all clauses. For all other classes, \CBa{} compile weak FOND problems into
deterministic ones (Lemma~4.3): every primitive task $a$ becomes a compound task $c_a$ with one
method per effect of $a$, whose network is a deterministic copy of $a$ with that single effect.
Choosing an outcome becomes choosing a method. The compilation preserves acyclicity, tail recursion
and total order (the new compound names form the lowest stratum), so weak plan existence has
exactly the complexity of deterministic plan existence in these classes
\parencite{erol:et:al:amai-96,alford:et:al:icaps-15}.

The compilation destroys regularity, since a method may now contain several compound tasks.
\CBa{} argue informally that the progression algorithm still works. A clean argument is
as follows. Regular problems are tail-recursive with a single stratum, and their compilation is
tail-recursive with two strata. Deterministic plan existence for tail-recursive problems of bounded
height is in \PSPACE, because progression keeps the network polynomially large
\parencite{alford:et:al:icaps-15}. Hence weak plan existence for regular problems is \PSPACE-complete.

\paragraph{Strong LD solutions.}
\CBa{} reduce strong LD plan existence to deterministic plan existence (Lemma~4.6) by replacing
every action by a deterministic action that adds what all effects add and deletes what any effect
deletes. The idea is that a single deterministic run tracks the intersection $\bigcap S_i$ of all
states the conformant prefix may lead to. With only positive preconditions this is correct: an
action is applicable in all states of $S_i$ iff its precondition is in $\bigcap S_i$. With
negative preconditions it is not, since a negative precondition must be false in \emph{every}
state of $S_i$, which the intersection cannot tell.

We therefore track both the intersection and the union, on two copies of the facts. An action is
applicable throughout $S$ iff its positive precondition is contained in $\bigcap S$ and its
negative precondition is disjoint from $\bigcup S$. Both sets evolve by ordinary STRIPS effects:
for the intersection, add what every effect adds and delete what some effect deletes without
adding it; for the union, add what some effect adds and delete what every effect deletes. Methods
and networks are unchanged, so every structural class is preserved, and strong LD plan existence
has the complexity of deterministic plan existence for every class. By
Proposition~\ref{prop:fm-relations}(iii) the same holds for strong OD solutions of totally ordered
problems \parencite[Cor.~4.8]{chen:bercher:icaps-21}.

The remaining entries of the fixed-method table concern strong OD solutions of partially ordered
problems, where policies may react to outcomes by reordering. This is where the complexity jumps,
and where most of our new results are.

\subsection{Primitive problems}

\begin{theorem}[\CBa, Thm.~5.1]\label{thm:prim-po}
  Strong OD plan existence for primitive partially ordered problems is \PSPACE-complete, even
  without negative preconditions and with at most two effects per action.
\end{theorem}

Membership is clear: by Proposition~\ref{prop:od-belief} we evaluate a game of depth $|T|$ on
configurations of polynomial size. For hardness, \CBa{} reduce from $(m,n,k)$-games
($k$-in-a-row on an $m\times n$ board). The construction does not work as stated. The first player
may place a stone on a square already occupied by the opponent, which triggers an ``illegal move''
task that unlocks a clean-up phase in which everything becomes executable. The task meant for the
first player's final move can be used at any time to make an unanswered extra move. And the task
recording a win re-enables the illegal-move check, so the final ``the opponent has no line'' test
can be bypassed. Each of these gives the planner a trivial win. Rather than repairing the gadget
we reduce from TQBF, the canonical \PSPACE-complete problem \parencite{stockmeyer:meyer:stoc-73}.

\begin{proof}[Proof of hardness]
  Let $\varphi=Q_1x_1\cdots Q_mx_m\,\psi$ with $\psi$ in conjunctive normal form. We use facts
  $x_i^0,x_i^1$ (``$x_i$ is false/true''), one fact $\mathit{sat}_j$ per clause, and a fact
  $\mathit{win}$. For each variable there is a \emph{block} of two tasks. If $x_i$ is existential,
  one task sets $x_i$ to true and the other sets it to false (each adds one literal fact and deletes
  the other). Both must be executed, and the one executed \emph{last} determines the value, so the
  policy chooses the value by choosing the order. If $x_i$ is universal, both tasks are the same
  nondeterministic action with outcomes ``set true'' and ``set false'', so the environment
  determines the value. All tasks of block $i$ precede all tasks of later blocks and all remaining
  tasks.

  After the blocks come \emph{check} tasks, one for each occurrence of a literal $\ell$ in a clause
  $j$, which require the fact of $\ell$ and add $\mathit{sat}_j$; a task $\mathit{fin}$, which
  requires all $\mathit{sat}_j$ and adds $\mathit{win}$; and a task $\mathit{clear}$, which
  requires $\mathit{win}$ and adds all literal facts. The check tasks, $\mathit{fin}$ and
  $\mathit{clear}$ are mutually unordered.

  While a block is open, only its two tasks are executable, and afterwards the state encodes the
  assignment chosen so far. After the last block, $\mathit{clear}$ needs $\mathit{win}$, and only
  $\mathit{fin}$ adds it. So the policy must make every clause satisfied using check tasks whose
  literals are true under the assignment, then run $\mathit{fin}$, then $\mathit{clear}$, and then
  the remaining check tasks (which $\mathit{clear}$ has made executable). This succeeds against
  every behaviour of the environment iff the existential player has a winning strategy in the QBF
  game, i.e.\ iff $\varphi$ is true. All preconditions are positive, and actions have at most two
  effects.
\end{proof}

By Proposition~\ref{prop:fm-vs-mp}(i) the same bound holds for strong and strong cyclic
method-based policies on primitive problems \parencite[Cor.~1]{chen:bercher:icaps-22}, whose
published proof inherits the flaw.

\subsection{Acyclic problems}\label{sec:fixed-acyc}

\begin{theorem}\label{thm:acyc-od}
  Strong OD plan existence for acyclic partially ordered problems is \EXPSPACE-complete.
\end{theorem}

The upper bound is due to \CBa{} (Thm.~5.2). In an acyclic problem every achievable network has
at most exponentially many tasks, and a nondeterministic exponential-space procedure can guess
the decompositions and then evaluate the game of Proposition~\ref{prop:od-belief}, which has
exponential depth and configurations of exponential size.

For the lower bound only \NEXPTIME-hardness was known, inherited from deterministic problems. An
outline in unpublished notes proposed to simulate a Turing machine with ``counter cells''. As
stated, it simulates a deterministic machine, and it leaves the gadgets open. We give a complete
construction based on the same idea, but with a different source problem.

\paragraph{The source: the Savitch game.}
Let $N$ be a nondeterministic Turing machine that uses $S=2^{p(n)}$ cells on inputs of length $n$.
Whether $N$ can move from a configuration $c$ to a configuration $c'$ in at most $2^k$ steps is
decided by the following game, which underlies Savitch's theorem. For $k=0$ the claim is checked
directly ($c=c'$ or $c'$ is a successor of $c$). For $k>0$ the existential player names a
midpoint $d$, and the universal player chooses which of the two claims $(c,d)$ and $(d,c')$, each
with $k-1$, to continue with. Configurations have exponential size and the game has exponentially
many rounds ($k$ is exponential in $n$), which is alternating exponential time and hence enough to
capture \EXPSPACE{} \parencite{chandra:et:al:jacm-81}. We use the formulation in which $w$ is
accepted iff, for some accepting configuration $e$, the existential player wins the claim
(initial configuration, $e$) with $k=2^{O(p(n))}$.

\paragraph{The memory: a queue of counter cells.}
A position of the game is a pair of configurations, i.e.\ two words of length $S$, too large for
the state of a planning problem. We store it in the \emph{task network}. The network consists of a
totally ordered \emph{control chain} and pools of mutually unordered \emph{cells}. A cell for the
symbol $\tau$ is a chain
\[
  \mathit{init}_\tau\cdot(\mathit{tick}\cdot\mathit{res})^S\cdot\mathit{fin}_\tau .
\]
The control chain is a long sequence of \emph{windows}. A window enables one kind of cell task (by
adding a token fact that the task consumes) and then disables it again. Each \emph{step} of the
control consists of $S$ windows for $\mathit{tick}$, $S$ windows for $\mathit{res}$, and then one
window each for $\mathit{fin}$, $\mathit{init}$ and an update of the finite control state. A cell
that has ticked must resume before it can tick again, so the $S$ tick windows of a step tick
\emph{every} running cell exactly once. Consequently the unique cell that can finish in a step is
the one started exactly $S$ steps earlier: $\mathit{fin}_\tau$ reports the symbol $\tau$ written $S$
steps before, and $\mathit{init}_{\tau'}$ lets the policy write a new symbol $\tau'$ of its choice
by starting a cell of that kind. The cells thus implement a \emph{queue} (a delay line) of length
$S$. Each step reads the symbol written $S$ steps earlier and appends a new one. A short
\emph{pre-warm} phase fills the queue with blanks. A policy that does not stagger the initial cells
perfectly gets stuck within $S$ steps.

\paragraph{The program.}
The finite control runs an automaton over the queue. In an input sweep of $S$ steps it writes, for
each cell $i$, the triple $(c_0(i),e(i),e(i))$ of the initial configuration and the claimed
accepting configuration $e$ (chosen by the policy). In each of the $k$ rounds it sweeps once more.
It reads the triple stored for cell $i$, projects it to the current pair $(x(i),y(i))$, and lets the
policy write $(x(i),m(i),y(i))$, where $m$ is the claimed midpoint. At the end of the round a
nondeterministic action chooses which half the next round projects to. After the last round a
check sweep verifies with a finite automaton that the remaining pair is equal or a single step of
$N$. When the check succeeds, a flag makes all remaining tasks executable. Since every symbol of
the current pair is read exactly $S$ steps after it was written, the policy wins iff it wins the
Savitch game.

\paragraph{Exponential repetition.}
The numbers $S$ and $k$, and the counts of windows and of cells, are exponential. They are produced
by acyclic \emph{doubling} hierarchies: a compound task $D_{j+1}$ has the single method
$D_{j+1}\to D_j\,D_j$, totally ordered for the control chain and unordered for the pools, and a
count given in binary selects which levels contribute. Every compound task has exactly one method,
so the whole problem decomposes into a single primitive network.

\begin{proof}[Proof of Theorem~\ref{thm:acyc-od}, hardness]
  The construction is computable in polynomial time, the problem is acyclic, and its unique
  achievable primitive network is the control chain together with the pools. By the queue property,
  every strong OD policy induces a strategy of the existential player in the Savitch game, which
  wins because the final check must succeed against every choice of halves. Conversely, a winning
  strategy is executed by writing the claimed midpoints. Hence $w$ is accepted iff the problem has
  a strong OD solution.
\end{proof}

Because every compound task has a single method, fixed methods and method-based policies coincide
on these instances. By Proposition~\ref{prop:fair-basic}(ii) strong and fair strong cyclic
policies coincide on acyclic problems. So the same reduction proves \EXPSPACE-hardness of strong
and strong cyclic method-based plan existence for acyclic partially ordered problems. This replaces
the proof of \textcite[Thm.~2]{chen:bercher:icaps-22}, which only outlines an extension of the
\NEXPTIME-hardness proof of \textcite{alford:et:al:icaps-15} to alternating machines.

\subsection{Regular problems}\label{sec:fixed-reg}

For regular problems the answer depends on the initial task network.

\begin{theorem}\label{thm:regtn}
  Strong OD plan existence for \REGtn{} is \EXPSPACE-complete.
\end{theorem}

The upper bound is due to \CBa{} (Thm.~5.3), whose argument we make precise. In a \REGtn{} problem
every achievable primitive network is a sequence of \emph{blocks} $B_0\prec B_1\prec\cdots\prec
B_\ell$: the primitive part of $\tnI$, then the primitive parts of the chosen methods, in order.
Every task of a block precedes every task of the next, because the compound task of a regular
network comes last and decomposition inherits its predecessors. The key fact is a
\emph{composition lemma}: $X$ followed by $Y$ is solvable from $s$ iff $X$ is solvable from $s$
in such a way that every final state lies in the set of states from which $Y$ is solvable. Hence
a solution is described by the sequence of methods and the sequence of \emph{belief sets}, the sets
of states in which the execution can be at the start of each block. A nondeterministic procedure
guesses the next method and the next belief set (an explicit set of states, of exponential size)
and checks the current block from each state of the current belief set in polynomial space. By
Savitch's theorem this is in \EXPSPACE.

\DB{} proved the matching lower bound via chosen-length planning games. Their proof leaves the
STRIPS gadgets of the game unspecified, and the simulation of games by regular problems is argued
only informally. We give a direct proof that is considerably simpler.

\begin{proof}[Proof of Theorem~\ref{thm:regtn}, hardness]
  Let $N$ be a nondeterministic machine with space $2^n$ and input $w$. The initial network starts
  with a nondeterministic task by which the environment chooses a cell index $c<2^n$, written in
  $n$ bits of the state, followed by a gadget that checks that a register $\sigma$ holds the input
  symbol at cell $c$, and then a single compound task $\mathit{run}$. The methods for
  $\mathit{run}$ are regular. For each transition of $N$ and each way of updating the binary head
  position there is a method $\mathit{run}\to B\cdot\mathit{run}$, and for each accepting state a
  final method $\mathit{run}\to B_{\mathit{acc}}$. The block $B$ of a transition from $q$ to $q'$
  that reads $a$ and writes $a'$ checks: if $c$ is the head position, then $\sigma=a$ and the new
  content is $a'$; otherwise the content stays $\sigma$. It then updates $\sigma$, the control
  state and the head position. The policy supplies the new content and auxiliary witnesses, which
  small CNF gadgets check.

  With fixed methods, the sequence of methods, i.e.\ a run of $N$, is chosen before the environment
  picks $c$, and it is the same in every branch. In the branch for $c$ the blocks track the
  content of cell $c$ along this run and fail iff the run is inconsistent at $c$. So a strong OD
  solution exists iff some accepting run of $N$ on $w$ is consistent at every cell, i.e.\ iff $N$
  accepts $w$. Every Boolean register is represented by a pair of facts, so only positive
  preconditions are used.
\end{proof}

The argument shows nicely why fixed methods are expensive here: the planner must commit to the
whole run before learning which cell is checked. With method-based policies the planner could
adapt the run to $c$, and indeed regular problems are only \EXPTIME-complete in that setting
(Section~\ref{sec:policies}).

For the class \REG{} of \textcite{bercher:et:al:ijcai-19}, where only methods must be regular, the
membership argument breaks down. With $\tnI$ consisting of several unordered compound tasks, the
achievable networks consist of several parallel block sequences. A policy may interleave them
differently in different branches, so there is no single cut between blocks at which a belief set
summarises the past. In fact no algorithm exists.

\begin{theorem}\label{thm:reg-undec}
  Strong OD plan existence is undecidable for problems that are both in \REG{} and in \TR{}. In
  particular it is undecidable for \REG{} and for partially ordered tail-recursive problems.
\end{theorem}

The undecidability of the tail-recursive case was conjectured in the same unpublished notes,
with a reduction from the Post correspondence problem that is to check two different equalities in
two ``modes'' of execution. The missing piece is a way to make the check of the inactive mode
free. With positive preconditions and unconditional effects this is not possible in the proposed
form, since the tasks of the inactive check still have to be executed. We solve this with
\emph{gated cheat tasks} and reduce directly from the halting problem of two-stack machines.

\begin{proof}
  Fix a two-stack machine $M$ and write a configuration with control state $q$ and stacks $u,v$
  as the word $\mathrm{rev}(u)\,q\,v$. A transition rewrites a factor of length two around $q$.
  We use letters that carry one optional character for each of two \emph{modes},
  written $(o_1,o_2)$.

  The initial network consists of four unordered compound tasks, each generating a chain by regular
  methods. Chain $A$ is $\mathit{init}$, followed by the letters $(\cdot,x)$ for the initial
  configuration and a separator $\#$, then arbitrary letters $(x,x)$ chosen by the decomposition,
  then the letters $(x,\cdot)$ for the halting configuration followed by $\#$, and finally a task
  $\mathit{fin}$. Chain $B$ is a sequence of blocks $u\cdot R(t)\cdot v\cdot\#$, where $u,v$ consist
  of letters $(x,x)$ and $R(t)$ lists the right-hand side of transition $t$ in mode~1 and its
  left-hand side in mode~2, followed by a task $\mathit{end}_B$. Projecting to mode~1 gives
  $\pi_1(A)=\mathit{mid}\cdot C_{\mathit{halt}}\#$ and $\pi_1(B)=$ the concatenation of the
  \emph{successor} configurations of the blocks. Projecting to mode~2 gives $\pi_2(A)=C_0\#\cdot
  \mathit{mid}$ and $\pi_2(B)=$ the concatenation of the \emph{predecessor} configurations.
  Both equalities $\pi_1(A)=\pi_1(B)$ and $\pi_2(A)=\pi_2(B)$ hold iff the blocks of $B$ describe a
  run $C_0\to C_1\to\cdots\to C_{\mathit{halt}}$ and $\mathit{mid}=C_1\#\cdots C_{k-1}\#$.

  The first task $\mathit{init}$ is nondeterministic and chooses the active mode $m$. A letter of
  $A$ with character $x$ in mode $m$ requires a token $\mathit{free}_m$ and replaces it by
  $\mathit{pend}_m(x)$, and a letter of $B$ with character $x$ in mode $m$ consumes
  $\mathit{pend}_m(x)$ and restores $\mathit{free}_m$. So in the active mode the letters of $A$ and
  $B$ must alternate character by character, which succeeds iff the projections agree. The task
  $\mathit{fin}$ requires $\mathit{free}_1$, $\mathit{free}_2$ and the fact added by
  $\mathit{end}_B$. The two remaining chains are \emph{supply chains} of cheat tasks
  $\mathit{cheat}_m$, whose length is chosen by the decomposition. A cheat adds all facts of mode
  $m$, so it makes every letter of mode $m$ executable, but it requires $\mathit{ok}_m$. The task
  $\mathit{init}$ adds $\mathit{ok}$ for the \emph{inactive} mode only, and $\mathit{fin}$ adds both.
  Thus in the inactive mode the policy can use cheats freely, while in the active mode cheats are
  available only after the check has been completed. The cheats sit in chains of their own
  because a cheat inside chain $A$ or $B$ would block that chain in the active mode.

  With fixed methods the words in $A$ and $B$ are chosen before $\mathit{init}$ reveals the mode,
  so a strong OD solution exists iff one choice passes both checks, i.e.\ iff $M$ halts. All
  methods are regular and the problem has a single stratum, so it lies in $\REG\cap\TR$.
\end{proof}

Strong OD plan existence is semi-decidable for all problems: enumerate the achievable primitive
networks and decide each in polynomial space \parencite[Thm.~3.1]{chen:bercher:icaps-21}. So
tail-recursive and \REG{} problems are as hard as unrestricted ones. For unrestricted partially
ordered problems all four fixed-method concepts are semi-decidable and undecidable, the latter
already for deterministic problems \parencite{erol:et:al:amai-96,geier:bercher:ijcai-11}.
